\begin{claim}{B.37}{Proved}{Complete genus-three mirror-disc quotient catalogue}
There are exactly $145$ conjugacy classes of finite subgroups $K$ of
$\mathcal T_{3,0}$ for which $H=K\cap\mathcal T^+_{3,0}$ has spherical
quotient and the induced involution on $S_3/H$ has a fixed circle.
Exactly $57$ contain a hyperelliptic involution, and $88$ do not.
The files \path{type_3_0_reflection_quotient_catalogue.csv},
\path{type_3_0_reflection_quotient_certificates.json}, and
\path{type_3_0_reflection_quotient_sections.json} give every
class, its exact finite group, and all its section data.
\end{claim}
\begin{proof}
The quotient by the induced involution is a mirror disc $Q$.  Let
$(m_1,\ldots,m_k)$ be its cyclic list of corner orders and
$(n_1,\ldots,n_l)$ its list of interior cone orders.  The orientation
double of $Q$ is the spherical orbifold with cone multiset
\[
               \{m_1,\ldots,m_k,n_1,n_1,\ldots,n_l,n_l\}.
                                                        \tag{B.37.1}
\]
The involution on a sphere with a fixed circle is conjugate to a
reflection: cut along that circle into two discs and extend a boundary
parametrization over either disc, then use the involution to define
the parametrization on the other disc.  Cone points on its circle
become the corners of $Q$, while its other cone points occur in
exchanged pairs.  Thus the corner and interior data in (B.37.1)
exhaust its quotient types.  The corner list is considered up to
cyclic permutation and reversal; the interior list is unordered.

Use the following presentation of $\Lambda=\pi_1^{\rm orb}(Q)$:
\[
 \left\langle x_1,\ldots,x_l,e,c_0,\ldots,c_k\ \left|
 \begin{array}{c}
 x_i^{n_i}=1,\quad x_1\cdots x_l e=1,\\
 c_j^2=1,\quad(c_{j-1}c_j)^{m_j}=1,\quad c_k=ec_0e^{-1}
 \end{array}\right.\right\rangle .                    \tag{B.37.2}
\]
For $k=0$ there is one reflection $c_0$, commuting with $e$.
The orientation kernel $\Delta=\Lambda^+$ is the orbifold group
of (B.37.1).  Put $t=c_0$, $Z_j=c_{j-1}c_j$,
$P_j=Z_1\cdots Z_j$, $U_i=x_{l+1-i}^{-1}$, and $V_i=tU_it$.
Schreier rewriting, or the doubled fundamental polygon, gives the
spherical generating vector
\[
 (Z_1,\ldots,Z_k,U_1,\ldots,U_l,V_l^{-1},\ldots,V_1^{-1}),
 \qquad P_k\,U_1\cdots U_l\,(V_1\cdots V_l)^{-1}=1.
                                                        \tag{B.37.3}
\]
For example $e=U_1\cdots U_l$ and
$P_k=t c_k=tet e^{-1}=(V_1\cdots V_l)e^{-1}$, which proves
the product relation with its stated order.  Conjugation by $t$
acts on these generators by
\[
 \psi(Z_j)=P_{j-1}Z_j^{-1}P_{j-1}^{-1},\qquad
 \psi(U_i)=V_i,\qquad\psi(V_i)=U_i.                    \tag{B.37.4}
\]
These formulas include every base-path conjugation.

The input \cite{R-KNSData}, whose completeness is established in
\cite[Theorem 2.2, Corollary 2.3, and the planar calculation]{R-KNS},
contains every $\operatorname{Aut}(H)$-class of smooth epimorphisms
$\theta:\Delta\to H$ in genus three.  Its $197$ vectors cover all
$49$ genus-three signature/group pairs, including those with positive
quotient genus.  Here we use precisely the spherical-quotient pairs.
The original input and its attribution are retained in
\texttt{research\_data/actions-genus3.json}; B.33 verifies its orders
and surface signatures.  A prescribed ordering (B.37.3) is obtained
from the published sorted ordering by explicit adjacent Hurwitz
moves.  This is a bijection of epimorphism sets, since each move has
its displayed inverse in the program.

For each vector the program tests whether its kernel is invariant
under (B.37.4).  More explicitly, it attempts to define
\[
              \alpha(\theta(a))=\theta(\psi(a))
                         \qquad(a\in\Delta)             \tag{B.37.5}
\]
by propagation on the finite Cayley graph.  Agreement on every edge
is exactly the condition that (B.37.5) be well defined.  Surjectivity
then makes it an automorphism and (B.37.4) gives $\alpha^2=1$.
The resulting full action group is explicitly
\[
           K=H\rtimes_\alpha\langle t\rangle,
                 \quad t^2=1.                         \tag{B.37.6}
\]
Its epimorphism from (B.37.2) sends $x_i$ to their recovered vector
entries and $c_j$ to $t\theta(P_j)$.  Its kernel is exactly
$\ker\theta\subset\Delta$.  In particular no choice of a further
extension class is omitted: the orbifold kernel determines
$\Lambda/\ker\theta$ and hence (B.37.6).  Conversely every action
in the claim supplies this invariant kernel and therefore occurs.

We now prove the precise fusion calculation.  Two epimorphisms from
the same orbifold group have the same kernel exactly when a unique
isomorphism between their finite images intertwines them: define
that isomorphism by $\theta(a)\mapsto\theta'(a)$, which is well defined
exactly when the kernels agree.  Thus the published
$\operatorname{Aut}(H)$-normalization records kernels without loss.
A conjugacy of the surface actions descends to a decorated-orbifold
homeomorphism of $Q$ carrying one kernel to the other.  Conversely
such a quotient homeomorphism lifts through the two regular
orbifold coverings because it carries their defining kernels to
one another; its lift is the required surface conjugacy.  This is
an exact equivalence, so it remains to calculate the full action of
these disc homeomorphisms on the finite kernel set.

A disc homeomorphism permutes its colored boundary corners cyclically,
or reverses their cyclic order.  After composing with that boundary
map it fixes the boundary pointwise up to isotopy.  The remaining
mapping class is a braid of the interior cone points.  To see this,
forget the interior points and isotope the boundary-fixed disc map
to the identity by the Alexander isotopy; the tracks of the points
are its braid, and an isotopy of those tracks gives precisely an
isotopy of the punctured disc map by isotopy extension.  The colored
braid subgroup is generated by pure braid generators and adjacent
exchanges among equal colors after the colors are sorted.  The
program uses the standard words
$\sigma_{j-1}\cdots\sigma_{i+1}\sigma_i^2
\sigma_{i+1}^{-1}\cdots\sigma_{j-1}^{-1}$ and those equal-color
exchanges.  On interior generators an adjacent positive move is
\[
                    (x_i,x_{i+1})\longmapsto
                    (x_ix_{i+1}x_i^{-1},x_i),           \tag{B.37.7}
\]
fixing $e$ and every boundary reflection.  The inverse is also used
when forming the pure braid words.

Here are the boundary moves, including their effect on the chosen
reflection.  Extend the sequence by $c_{j+k}=ec_je^{-1}$.  One cyclic
shift has
\[
 \begin{split}
 (Z_1,\ldots,Z_k)&\longmapsto
       (Z_2,\ldots,Z_k,eZ_1e^{-1}),\\
 (x_1,\ldots,x_l)&\longmapsto(x_1,\ldots,x_l),\\
 t&\longmapsto tZ_1,\qquad
 \alpha\longmapsto\operatorname{Inn}(Z_1^{-1})\alpha.
 \end{split}                                           \tag{B.37.8}
\]
A reversal has
\[
 \begin{split}
 (Z_1,\ldots,Z_k)&\longmapsto(Z_k^{-1},\ldots,Z_1^{-1}),\\
 (x_1,\ldots,x_l)&\longmapsto(x_l^{-1},\ldots,x_1^{-1}),\\
 t&\longmapsto tP_k,\qquad
 \alpha\longmapsto\operatorname{Inn}(P_k^{-1})\alpha.
 \end{split}                                           \tag{B.37.9}
\]
They follow by sending $c_j$ respectively to $c_{j+1}$ and
$c_{k-j}$ in (B.37.2); in the reversal $e$ goes to $e^{-1}$.
Only the shifts and shifted reversals preserving the corner-order
list are self-homeomorphisms of the chosen colored quotient.
A fixed braid restores the sorted interior order after a reversal.
Equations (B.37.7)--(B.37.9), together with the preceding braid
argument, prove both that every programmed fusion is realized and
that every actual conjugacy is included.

The program computes the finite orbits under these moves.  It
identifies two generating vectors under $\operatorname{Aut}(H)$ by
their rooted, edge-colored Cayley graphs: starting at the identity,
visit vertices in generator order and record each colored successor.
Equality defines a unique group isomorphism mapping each generator
to the corresponding generator, and such an isomorphism plainly
gives equality.  Thus this is an exact isomorphism test rather than
a coarse group invariant.  The computation gives $145$ orbits.

For each orbit the second program constructs the regular action of
(B.37.6) explicitly from
$(h,\epsilon)(h',\epsilon')=(h\alpha^\epsilon(h'),\epsilon+\epsilon')$.
It checks (B.37.2), every local stabilizer injection, and
$|K|\chi^{\rm orb}(Q)=-4$.  The resulting cover is therefore a closed
orientable genus-three surface.  The finite coset incidence graph
then gives every fixed point and oval, multiplier, section,
extension, maximal section, centralizer, orbit and stabilizer as in
B.26.  These checks passed for all $145$ classes.

Finally the overlap with B.34 is identified by actual monodromy
isomorphisms.  In each of B.34's $57$ mirror-disc groups the program
extracts the positive group and the generators $x_i$, $Z_j$ and
conjugation by $c_0$.  It applies (B.37.8)--(B.37.9) and the interior
braids to put the signature in the same ordering, and locates its
rooted colored Cayley graph in the computed kernel orbit.  The file
\texttt{genus3\_hyperelliptic\_reflection\_crosscheck.json} records
the resulting bijection with $57$ different rows and equality of
all eight reported section-count columns.  Thus these are exactly
the common classes, leaving $88$ additional ones.
\end{proof}

The $145$ classes have the following order distribution:
\[
\begin{array}{c|rrrrrrrr}
 |K|&4&6&8&12&14&16&18&24\\\hline
 \#&5&3&36&13&3&35&1&15\\[2pt]
 |K|&28&32&42&48&64&96&192&336\\\hline
 \#&1&13&1&11&2&4&1&1.
\end{array}                                             \tag{B.37.10}
\]
Together B.34 and B.37 therefore give $149$ distinct mixed classes.
The additional four are the projective-plane quotients in B.34.
The scope of B.37 requires both a spherical positive quotient and a
fixed-circle involution on it; the remaining quotient types are
listed in the accompanying computation report.
